% Options for packages loaded elsewhere
\PassOptionsToPackage{unicode}{hyperref}
\PassOptionsToPackage{hyphens}{url}
\documentclass[
]{article}
\usepackage{xcolor}
\usepackage{amsmath,amssymb}
\setcounter{secnumdepth}{-\maxdimen} % remove section numbering
\usepackage{iftex}
\ifPDFTeX
  \usepackage[T1]{fontenc}
  \usepackage[utf8]{inputenc}
  \usepackage{textcomp} % provide euro and other symbols
\else % if luatex or xetex
  \usepackage{unicode-math} % this also loads fontspec
  \defaultfontfeatures{Scale=MatchLowercase}
  \defaultfontfeatures[\rmfamily]{Ligatures=TeX,Scale=1}
\fi
\usepackage{lmodern}
\ifPDFTeX\else
  % xetex/luatex font selection
\fi
% Use upquote if available, for straight quotes in verbatim environments
\IfFileExists{upquote.sty}{\usepackage{upquote}}{}
\IfFileExists{microtype.sty}{% use microtype if available
  \usepackage[]{microtype}
  \UseMicrotypeSet[protrusion]{basicmath} % disable protrusion for tt fonts
}{}
\makeatletter
\@ifundefined{KOMAClassName}{% if non-KOMA class
  \IfFileExists{parskip.sty}{%
    \usepackage{parskip}
  }{% else
    \setlength{\parindent}{0pt}
    \setlength{\parskip}{6pt plus 2pt minus 1pt}}
}{% if KOMA class
  \KOMAoptions{parskip=half}}
\makeatother
\usepackage{color}
\usepackage{fancyvrb}
\newcommand{\VerbBar}{|}
\newcommand{\VERB}{\Verb[commandchars=\\\{\}]}
\DefineVerbatimEnvironment{Highlighting}{Verbatim}{commandchars=\\\{\}}
% Add ',fontsize=\small' for more characters per line
\newenvironment{Shaded}{}{}
\newcommand{\AlertTok}[1]{\textcolor[rgb]{1.00,0.00,0.00}{\textbf{#1}}}
\newcommand{\AnnotationTok}[1]{\textcolor[rgb]{0.38,0.63,0.69}{\textbf{\textit{#1}}}}
\newcommand{\AttributeTok}[1]{\textcolor[rgb]{0.49,0.56,0.16}{#1}}
\newcommand{\BaseNTok}[1]{\textcolor[rgb]{0.25,0.63,0.44}{#1}}
\newcommand{\BuiltInTok}[1]{\textcolor[rgb]{0.00,0.50,0.00}{#1}}
\newcommand{\CharTok}[1]{\textcolor[rgb]{0.25,0.44,0.63}{#1}}
\newcommand{\CommentTok}[1]{\textcolor[rgb]{0.38,0.63,0.69}{\textit{#1}}}
\newcommand{\CommentVarTok}[1]{\textcolor[rgb]{0.38,0.63,0.69}{\textbf{\textit{#1}}}}
\newcommand{\ConstantTok}[1]{\textcolor[rgb]{0.53,0.00,0.00}{#1}}
\newcommand{\ControlFlowTok}[1]{\textcolor[rgb]{0.00,0.44,0.13}{\textbf{#1}}}
\newcommand{\DataTypeTok}[1]{\textcolor[rgb]{0.56,0.13,0.00}{#1}}
\newcommand{\DecValTok}[1]{\textcolor[rgb]{0.25,0.63,0.44}{#1}}
\newcommand{\DocumentationTok}[1]{\textcolor[rgb]{0.73,0.13,0.13}{\textit{#1}}}
\newcommand{\ErrorTok}[1]{\textcolor[rgb]{1.00,0.00,0.00}{\textbf{#1}}}
\newcommand{\ExtensionTok}[1]{#1}
\newcommand{\FloatTok}[1]{\textcolor[rgb]{0.25,0.63,0.44}{#1}}
\newcommand{\FunctionTok}[1]{\textcolor[rgb]{0.02,0.16,0.49}{#1}}
\newcommand{\ImportTok}[1]{\textcolor[rgb]{0.00,0.50,0.00}{\textbf{#1}}}
\newcommand{\InformationTok}[1]{\textcolor[rgb]{0.38,0.63,0.69}{\textbf{\textit{#1}}}}
\newcommand{\KeywordTok}[1]{\textcolor[rgb]{0.00,0.44,0.13}{\textbf{#1}}}
\newcommand{\NormalTok}[1]{#1}
\newcommand{\OperatorTok}[1]{\textcolor[rgb]{0.40,0.40,0.40}{#1}}
\newcommand{\OtherTok}[1]{\textcolor[rgb]{0.00,0.44,0.13}{#1}}
\newcommand{\PreprocessorTok}[1]{\textcolor[rgb]{0.74,0.48,0.00}{#1}}
\newcommand{\RegionMarkerTok}[1]{#1}
\newcommand{\SpecialCharTok}[1]{\textcolor[rgb]{0.25,0.44,0.63}{#1}}
\newcommand{\SpecialStringTok}[1]{\textcolor[rgb]{0.73,0.40,0.53}{#1}}
\newcommand{\StringTok}[1]{\textcolor[rgb]{0.25,0.44,0.63}{#1}}
\newcommand{\VariableTok}[1]{\textcolor[rgb]{0.10,0.09,0.49}{#1}}
\newcommand{\VerbatimStringTok}[1]{\textcolor[rgb]{0.25,0.44,0.63}{#1}}
\newcommand{\WarningTok}[1]{\textcolor[rgb]{0.38,0.63,0.69}{\textbf{\textit{#1}}}}
\usepackage{longtable,booktabs,array}
\usepackage{calc} % for calculating minipage widths
% Correct order of tables after \paragraph or \subparagraph
\usepackage{etoolbox}
\makeatletter
\patchcmd\longtable{\par}{\if@noskipsec\mbox{}\fi\par}{}{}
\makeatother
% Allow footnotes in longtable head/foot
\IfFileExists{footnotehyper.sty}{\usepackage{footnotehyper}}{\usepackage{footnote}}
\makesavenoteenv{longtable}
\usepackage{graphicx}
\makeatletter
\newsavebox\pandoc@box
\newcommand*\pandocbounded[1]{% scales image to fit in text height/width
  \sbox\pandoc@box{#1}%
  \Gscale@div\@tempa{\textheight}{\dimexpr\ht\pandoc@box+\dp\pandoc@box\relax}%
  \Gscale@div\@tempb{\linewidth}{\wd\pandoc@box}%
  \ifdim\@tempb\p@<\@tempa\p@\let\@tempa\@tempb\fi% select the smaller of both
  \ifdim\@tempa\p@<\p@\scalebox{\@tempa}{\usebox\pandoc@box}%
  \else\usebox{\pandoc@box}%
  \fi%
}
% Set default figure placement to htbp
\def\fps@figure{htbp}
\makeatother
\setlength{\emergencystretch}{3em} % prevent overfull lines
\providecommand{\tightlist}{%
  \setlength{\itemsep}{0pt}\setlength{\parskip}{0pt}}
\usepackage{multicol}
\usepackage{bookmark}
\IfFileExists{xurl.sty}{\usepackage{xurl}}{} % add URL line breaks if available
\urlstyle{same}
\hypersetup{
  pdftitle={pandoc-utils Feature Demo},
  hidelinks,
  pdfcreator={LaTeX via pandoc}}

\title{pandoc-utils Feature Demo}
\author{}
\date{}

\begin{document}
\maketitle

\section{CSV Tables}\label{csv-tables}

\begin{longtable}[]{@{}ccc@{}}
\caption{Benchmark Results}\tabularnewline
\toprule\noalign{}
n & time\_ms & error \\
\midrule\noalign{}
\endfirsthead
\toprule\noalign{}
n & time\_ms & error \\
\midrule\noalign{}
\endhead
\bottomrule\noalign{}
\endlastfoot
10.0 & 0.0 & 9.47e-17 \\
100.0 & 0.44 & 5.9e-16 \\
1000.0 & 3.75 & 1.48e-15 \\
10000.0 & 39.9 & 8.27e-16 \\
\end{longtable}

\section{Mermaid Diagrams}\label{mermaid-diagrams}

\pandocbounded{\includegraphics[keepaspectratio]{mermaid-images/eb89df5389d37301b76912c9a36e6c9216078328.pdf}}

\pandocbounded{\includegraphics[keepaspectratio]{mermaid-images/6b7ed204e608e78bdb28298f700be677fc4b96f9.pdf}}

\section{SVG Images}\label{svg-images}

\begin{figure}
\centering
\pandocbounded{\includegraphics[keepaspectratio,alt={Pandoc pipeline}]{demo_pipeline.pdf}}
\caption{Pandoc pipeline}
\end{figure}

\section{Include Files}\label{include-files}

\subsection{pandoc-utils}\label{pandoc-utils}

Utilities for my pandoc/markdown/LaTeX workflows.

\subsubsection{Structure}\label{structure}

\begin{verbatim}
pandoc-utils/
|-- pandoc_md_pdf.sh                # markdown -> LaTeX -> PDF pipeline
|-- Makefile                        # dep-install, dep-check, test, demo targets
|-- pyproject.toml                  # Python deps managed via uv
|-- markdown_get_defaults_yaml.py   # (internal) extract pandoc config from YAML front matter
|-- filters/
|   |-- svg_to_pdf.py               # convert SVG images to PDF for LaTeX
|   |-- mermaid.py                  # render mermaid diagrams to SVG/PDF
|   |-- csv_tex_table.py            # CSV code blocks to tables
|   |-- include-files.lua           # include external markdown files
|   |-- include-code-files.lua      # include source code snippets
|   |-- columns.lua                 # multi-column layout (HTML + LaTeX)
|   +-- table.lua                   # YAML-driven table templating
+-- html-papers/                    # a markdown paper as a web page and a PDF (see its README)
    |-- html-papers.lua             # the filter: numbered references, terms, subfigures, ...
    |-- demo/                       # a paper that builds itself: paper.md, build.sh, its experiment
    |-- html-papers.css             # the layout: text, notes margin, contents, comments
    |-- paper.css                   # the look: title block, tables, figures, bibliography
    |-- popups.js                   # margin notes and previews on hover
    |-- contents.js                 # contents sidebar
    |-- live-figures.js             # click a static figure for its live version
    |-- arxiv-links.js              # links in the bibliography
    |-- run-code.js                 # code blocks the reader can edit and run (Python in the browser)
    |-- keep-place.js               # keeps the reader's place when the text reflows
    |-- menu.js                     # the reader's settings (M)
    |-- review.js                   # comments and suggested edits, exported as JSON
    +-- collect_reviews.py          # merge reviewers' files, list comments by source line
\end{verbatim}

\subsubsection{Core Scripts}\label{core-scripts}

\paragraph{pandoc\_md\_pdf.sh}\label{pandoc_md_pdf.sh}

Main conversion pipeline: markdown -\textgreater{} LaTeX -\textgreater{}
PDF with optional bibliography support.

\begin{Shaded}
\begin{Highlighting}[]
\ExtensionTok{pandoc\_md\_pdf.sh}\NormalTok{ document.md       }\CommentTok{\# convert to PDF (biblatex auto{-}detected)}
\ExtensionTok{pandoc\_md\_pdf.sh}\NormalTok{ document.md }\AttributeTok{{-}b}    \CommentTok{\# force biblatex on}
\end{Highlighting}
\end{Shaded}

Extracts pandoc configuration from the document's YAML front matter (via
\texttt{uv\ run\ markdown\_get\_defaults\_yaml.py}), then runs pandoc
followed by \texttt{latexmk}. Biblatex is enabled automatically when the
front matter declares a \texttt{bibliography:} or \texttt{references:}
field (pass \texttt{-b} to force it on regardless); when it is off, no
bibliography markup is emitted so documents without citations build
cleanly. \texttt{latexmk} runs \texttt{pdflatex} and, when biblatex is
enabled, \texttt{biber}, re-running each as many times as needed until
references converge. Produces \texttt{.tex}, \texttt{.pdf}, and
\texttt{.pdlog} files.

Internal helpers

\begin{itemize}
\tightlist
\item
  \textbf{markdown\_get\_defaults\_yaml.py} --- Extracts pandoc defaults
  from YAML front matter and writes a defaults file. Called internally
  by \texttt{pandoc\_md\_pdf.sh}. Resolves the \texttt{\$FILTERS\$}
  placeholder to the actual filters directory.
\end{itemize}

\subsubsection{Filters}\label{filters}

Pandoc filters are referenced in YAML front matter using the
\texttt{\$FILTERS\$} placeholder:

\begin{Shaded}
\begin{Highlighting}[]
\PreprocessorTok{{-}{-}{-}}
\FunctionTok{\_\_defaults\_\_}\KeywordTok{:}
\AttributeTok{  }\FunctionTok{filters}\KeywordTok{:}
\AttributeTok{    }\KeywordTok{{-}}\AttributeTok{ $FILTERS$/mermaid.py}
\AttributeTok{    }\KeywordTok{{-}}\AttributeTok{ $FILTERS$/svg\_to\_pdf.py}
\AttributeTok{    }\KeywordTok{{-}}\AttributeTok{ $FILTERS$/columns.lua}
\PreprocessorTok{{-}{-}{-}}
\end{Highlighting}
\end{Shaded}

\paragraph{svg\_to\_pdf.py}\label{svg_to_pdf.py}

Converts \texttt{.svg} images to PDF using Inkscape when targeting LaTeX
output. Caches results --- only re-converts when the SVG is newer than
the existing PDF.

\paragraph{mermaid.py}\label{mermaid.py}

Renders \texttt{mermaid} code blocks to SVG (HTML output) or PDF (LaTeX
output) using \texttt{mmdc}. For LaTeX, converts SVG -\textgreater{} PDF
via Inkscape and crops with \texttt{pdfcrop}. Caches output by content
hash.

\paragraph{csv\_tex\_table.py}\label{csv_tex_table.py}

Converts \texttt{csvtable} code blocks containing CSV data into pandoc
tables. (WIP)

\paragraph{include-files.lua}\label{include-files.lua}

Includes external markdown files at compile time. Supports recursive
inclusion, heading level shifting, relative path resolution, and comment
lines.

\begin{Shaded}
\begin{Highlighting}[]
\InformationTok{\textasciigrave{}\textasciigrave{}\textasciigrave{}include}
\InformationTok{chapter1.md}
\InformationTok{chapter2.md}
\InformationTok{\textasciigrave{}\textasciigrave{}\textasciigrave{}}
\end{Highlighting}
\end{Shaded}

\paragraph{include-code-files.lua}\label{include-code-files.lua}

Includes source code from external files into code blocks, with optional
line range selection and dedentation.

\begin{Shaded}
\begin{Highlighting}[]
\InformationTok{\textasciigrave{}\textasciigrave{}\textasciigrave{}python \{include="src/main.py" startLine=10 endLine=25 dedent=4\}}
\InformationTok{\textasciigrave{}\textasciigrave{}\textasciigrave{}}
\end{Highlighting}
\end{Shaded}

\paragraph{columns.lua}\label{columns.lua}

Multi-column layout support for both HTML (CSS columns) and LaTeX
(\texttt{multicols}). Configurable column count, gap, rules, and
justification.

\begin{Shaded}
\begin{Highlighting}[]
\NormalTok{::: \{.columns column{-}count=3 column{-}gap="1em"\}}
\NormalTok{Column 1 content}

\NormalTok{::: columnbreak}
\NormalTok{:::}

\NormalTok{Column 2 content}
\NormalTok{:::}
\end{Highlighting}
\end{Shaded}

\paragraph{table.lua}\label{table.lua}

YAML-driven table generation with reusable templates, configurable
alignment and widths, auto-numbering, and horizontal orientation
support.

\subsubsection{HTML papers}\label{html-papers}

\href{html-papers/README.md}{\texttt{html-papers/}} builds a paper
written in markdown as a web page with three columns: the contents, the
text, and a side column with notes and previews (citations, footnotes,
terms, cross-references) and, in review mode, comments and suggested
edits. The columns resize, and on a phone the notes open in the text.
Figures can be live, and a reader can edit and run the paper's code on
the page. The same source builds the PDF. A filter, two stylesheets,
eight drop-in scripts, and a script that maps reviewers' comments back
to lines of the source.

The demo is \href{https://mivanit.github.io/pandoc-utils/html-papers/}{a
paper that builds itself}: its build runs its experiment and prints the
code that did.

\begin{Shaded}
\begin{Highlighting}[]
\FunctionTok{make}\NormalTok{ demo{-}html{-}papers        }\CommentTok{\# html{-}papers/demo/ {-}\textgreater{} docs/html{-}papers/index.html and paper.pdf}
\FunctionTok{make}\NormalTok{ demo{-}html{-}papers{-}html   }\CommentTok{\# the web page only, in a few seconds}
\end{Highlighting}
\end{Shaded}

\subsubsection{Setup}\label{setup}

Dependencies are managed via \texttt{Makefile} and
\texttt{pyproject.toml}:

\begin{Shaded}
\begin{Highlighting}[]
\FunctionTok{make}\NormalTok{ dep{-}install   }\CommentTok{\# install system (apt), npm, and Python (uv) deps}
\FunctionTok{make}\NormalTok{ dep{-}check     }\CommentTok{\# verify all required binaries are available}
\FunctionTok{make}\NormalTok{ test          }\CommentTok{\# run pytest}
\FunctionTok{make}\NormalTok{ test{-}html     }\CommentTok{\# html{-}papers/ in headless Chromium (once: uv run playwright install chromium)}
\FunctionTok{make}\NormalTok{ test{-}html{-}firefox  }\CommentTok{\# the same in Firefox (once: uv run playwright install firefox)}
\end{Highlighting}
\end{Shaded}

\paragraph{Dependencies}\label{dependencies}

\textbf{System binaries:} \texttt{uv}, \texttt{pandoc}
(\textgreater=2.12), \texttt{latexmk}, \texttt{pdflatex},
\texttt{inkscape}, \texttt{mmdc} (mermaid-cli), \texttt{pdfcrop},
\texttt{magick} (ImageMagick), \texttt{biber} (for bibliographies)

\textbf{Python} (managed by \texttt{uv} via \texttt{pyproject.toml}):
\texttt{pyyaml}, \texttt{pandocfilters}, \texttt{pandas}

\section{Include Code Files}\label{include-code-files}

Full file:

\begin{Shaded}
\begin{Highlighting}[]
\CommentTok{\# /// script}
\CommentTok{\# requires{-}python = "\textgreater{}=3.10"}
\CommentTok{\# dependencies = [}
\CommentTok{\#     "pandocfilters",}
\CommentTok{\#     "pandas",}
\CommentTok{\# ]}
\CommentTok{\# ///}

\CommentTok{"""}
\CommentTok{pandoc filter to process code blocks with class "csvtable",}
\CommentTok{turning csv{-}formatted data into a table}
\CommentTok{"""}

\ImportTok{from}\NormalTok{ typing }\ImportTok{import}\NormalTok{ Any}
\ImportTok{from}\NormalTok{ io }\ImportTok{import}\NormalTok{ StringIO}

\ImportTok{import}\NormalTok{ pandas }\ImportTok{as}\NormalTok{ pd}
\ImportTok{from}\NormalTok{ pandocfilters }\ImportTok{import}\NormalTok{ toJSONFilter, get\_caption}


\NormalTok{READ\_KWARGS\_DEFAULT: }\BuiltInTok{dict}\NormalTok{[}\BuiltInTok{str}\NormalTok{, Any] }\OperatorTok{=}\NormalTok{ \{\}}


\KeywordTok{def}\NormalTok{ \_emptyblock() }\OperatorTok{{-}\textgreater{}} \BuiltInTok{list}\NormalTok{:}
    \ControlFlowTok{return}\NormalTok{ [}\StringTok{""}\NormalTok{, [], []]}


\KeywordTok{def}\NormalTok{ \_plain\_factory(val: }\BuiltInTok{str}\NormalTok{) }\OperatorTok{{-}\textgreater{}} \BuiltInTok{dict}\NormalTok{:}
    \ControlFlowTok{return}\NormalTok{ \{}
        \StringTok{"t"}\NormalTok{: }\StringTok{"Plain"}\NormalTok{,}
        \StringTok{"c"}\NormalTok{: [\{}\StringTok{"t"}\NormalTok{: }\StringTok{"Str"}\NormalTok{, }\StringTok{"c"}\NormalTok{: val.strip()\}],}
\NormalTok{    \}}


\KeywordTok{def}\NormalTok{ \_table\_cell\_factory(val: }\BuiltInTok{str}\NormalTok{) }\OperatorTok{{-}\textgreater{}} \BuiltInTok{list}\NormalTok{:}
    \ControlFlowTok{return}\NormalTok{ [}
\NormalTok{        \_emptyblock(),}
\NormalTok{        \{}\StringTok{"t"}\NormalTok{: }\StringTok{"AlignDefault"}\NormalTok{\},}
        \DecValTok{1}\NormalTok{,}
        \DecValTok{1}\NormalTok{,}
\NormalTok{        [\_plain\_factory(val)],}
\NormalTok{    ]}


\KeywordTok{def}\NormalTok{ \_table\_row\_factory(vals: }\BuiltInTok{list}\NormalTok{[}\BuiltInTok{str}\NormalTok{]) }\OperatorTok{{-}\textgreater{}} \BuiltInTok{list}\NormalTok{:}
    \ControlFlowTok{return}\NormalTok{ [\_emptyblock(), [\_table\_cell\_factory(v) }\ControlFlowTok{for}\NormalTok{ v }\KeywordTok{in}\NormalTok{ vals]]}


\KeywordTok{def}\NormalTok{ \_header\_factory(vals: }\BuiltInTok{list}\NormalTok{[}\BuiltInTok{str}\NormalTok{]) }\OperatorTok{{-}\textgreater{}} \BuiltInTok{list}\NormalTok{:}
    \ControlFlowTok{return}\NormalTok{ [\_emptyblock(), [\_table\_row\_factory(vals)]]}


\KeywordTok{def}\NormalTok{ \_body\_factory(rows: }\BuiltInTok{list}\NormalTok{[}\BuiltInTok{list}\NormalTok{[}\BuiltInTok{str}\NormalTok{]]) }\OperatorTok{{-}\textgreater{}} \BuiltInTok{list}\NormalTok{:}
    \ControlFlowTok{return}\NormalTok{ [}
\NormalTok{        [}
\NormalTok{            \_emptyblock(),}
            \DecValTok{0}\NormalTok{,}
\NormalTok{            [],}
\NormalTok{            [\_table\_row\_factory(row) }\ControlFlowTok{for}\NormalTok{ row }\KeywordTok{in}\NormalTok{ rows],}
\NormalTok{        ]}
\NormalTok{    ]}


\KeywordTok{def}\NormalTok{ rawcsv\_to\_df(}
\NormalTok{    in\_csv: }\BuiltInTok{str}\NormalTok{,}
\NormalTok{    read\_kwargs: }\BuiltInTok{dict}\NormalTok{[}\BuiltInTok{str}\NormalTok{, Any] }\OperatorTok{|} \VariableTok{None} \OperatorTok{=} \VariableTok{None}\NormalTok{,}
\NormalTok{) }\OperatorTok{{-}\textgreater{}}\NormalTok{ pd.DataFrame:}
    \ControlFlowTok{if}\NormalTok{ read\_kwargs }\KeywordTok{is} \VariableTok{None}\NormalTok{:}
\NormalTok{        read\_kwargs }\OperatorTok{=}\NormalTok{ READ\_KWARGS\_DEFAULT}
    \ControlFlowTok{return}\NormalTok{ pd.read\_csv(StringIO(in\_csv), }\OperatorTok{**}\NormalTok{read\_kwargs)}


\KeywordTok{def}\NormalTok{ csv\_tex\_table(}
\NormalTok{    key: }\BuiltInTok{str}\NormalTok{, value: Any, }\BuiltInTok{format}\NormalTok{: }\BuiltInTok{str}\NormalTok{, meta: Any}
\NormalTok{) }\OperatorTok{{-}\textgreater{}} \BuiltInTok{dict}\NormalTok{[}\BuiltInTok{str}\NormalTok{, Any] }\OperatorTok{|} \VariableTok{None}\NormalTok{:}
    \ControlFlowTok{if}\NormalTok{ key }\OperatorTok{!=} \StringTok{"CodeBlock"}\NormalTok{:}
        \ControlFlowTok{return} \VariableTok{None}

\NormalTok{    [[ident, classes, keyvals], code] }\OperatorTok{=}\NormalTok{ value}

    \ControlFlowTok{if} \StringTok{"csvtable"} \KeywordTok{not} \KeywordTok{in}\NormalTok{ classes:}
        \ControlFlowTok{return} \VariableTok{None}

\NormalTok{    caption, typef, keyvals }\OperatorTok{=}\NormalTok{ get\_caption(keyvals)}
\NormalTok{    df }\OperatorTok{=}\NormalTok{ rawcsv\_to\_df(code)}

\NormalTok{    n\_cols }\OperatorTok{=} \BuiltInTok{len}\NormalTok{(df.columns)}
\NormalTok{    colspecs }\OperatorTok{=}\NormalTok{ [[\{}\StringTok{"t"}\NormalTok{: }\StringTok{"AlignCenter"}\NormalTok{\}, \{}\StringTok{"t"}\NormalTok{: }\StringTok{"ColWidthDefault"}\NormalTok{\}] }\ControlFlowTok{for}\NormalTok{ \_ }\KeywordTok{in} \BuiltInTok{range}\NormalTok{(n\_cols)]}
\NormalTok{    header }\OperatorTok{=}\NormalTok{ \_header\_factory([}\BuiltInTok{str}\NormalTok{(c) }\ControlFlowTok{for}\NormalTok{ c }\KeywordTok{in}\NormalTok{ df.columns])}
\NormalTok{    body }\OperatorTok{=}\NormalTok{ \_body\_factory([[}\BuiltInTok{str}\NormalTok{(v) }\ControlFlowTok{for}\NormalTok{ v }\KeywordTok{in}\NormalTok{ row] }\ControlFlowTok{for}\NormalTok{ row }\KeywordTok{in}\NormalTok{ df.to\_numpy().tolist()])}
\NormalTok{    foot }\OperatorTok{=}\NormalTok{ [\_emptyblock(), []]}
\NormalTok{    cap }\OperatorTok{=}\NormalTok{ [}
        \VariableTok{None}\NormalTok{,}
\NormalTok{        []}
        \ControlFlowTok{if} \KeywordTok{not}\NormalTok{ caption}
        \ControlFlowTok{else}\NormalTok{ [\_plain\_factory(}\StringTok{" "}\NormalTok{.join(s[}\StringTok{"c"}\NormalTok{] }\ControlFlowTok{for}\NormalTok{ s }\KeywordTok{in}\NormalTok{ caption }\ControlFlowTok{if} \StringTok{"c"} \KeywordTok{in}\NormalTok{ s))],}
\NormalTok{    ]}

    \ControlFlowTok{return}\NormalTok{ \{}
        \StringTok{"t"}\NormalTok{: }\StringTok{"Table"}\NormalTok{,}
        \StringTok{"c"}\NormalTok{: [\_emptyblock(), cap, colspecs, header, body, foot],}
\NormalTok{    \}}


\KeywordTok{def}\NormalTok{ main() }\OperatorTok{{-}\textgreater{}} \VariableTok{None}\NormalTok{:}
    \CommentTok{"""Entry point for the pandoc filter."""}
\NormalTok{    toJSONFilter(csv\_tex\_table)}


\ControlFlowTok{if} \VariableTok{\_\_name\_\_} \OperatorTok{==} \StringTok{"\_\_main\_\_"}\NormalTok{:}
\NormalTok{    main()}
\end{Highlighting}
\end{Shaded}

Line range with dedent:

\begin{Shaded}
\begin{Highlighting}[]
\CommentTok{{-}{-}[[{-}{-} \# Columns {-} multiple column support in Pandoc\textquotesingle{}s markdown.}

\CommentTok{This Lua filter provides support for multiple columns in}
\CommentTok{latex and html outputs. For details, see README.md.}

\CommentTok{from: https://github.com/dialoa/columns/blob/master/columns.lua}

\CommentTok{@author Julien Dutant \textless{}julien.dutant@kcl.ac.uk\textgreater{}}
\CommentTok{@copyright 2021 Julien Dutant}
\CommentTok{@license MIT {-} see LICENSE file for details.}
\end{Highlighting}
\end{Shaded}

\section{Multi-Column Layout}\label{multi-column-layout}

{\setlength{\columnsep}{1.5em}
\begin{multicols}{3}

\textbf{Column 1} --- Filters process the document AST, transforming
custom code blocks and divs into native pandoc elements.

\columnbreak

\textbf{Column 2} --- Output adapts to the target format: LaTeX uses
\texttt{multicols}, HTML uses CSS columns.

\columnbreak

\textbf{Column 3} --- All filters can be combined freely in a single
document via the YAML front matter.

\end{multicols}
}

\subsection{Shorthand syntax}\label{shorthand-syntax}

{\begin{multicols}{2}

Left column content goes here. The \texttt{two-columns} class is
syntactic sugar for \texttt{\{.columns\ column-count=2\}}.

\columnbreak

Right column content goes here.

\end{multicols}
}

\subsection{Column span}\label{column-span}

{\begin{multicols}{2}

Left column.

\columnbreak

Right column.

\end{multicols}
}

This paragraph spans across all columns.

\section{YAML Tables}\label{yaml-tables}

\subsection{Vertical table with
auto-numbering}\label{vertical-table-with-auto-numbering}

\protect\phantomsection\label{tbl:team}
\begin{longtable}[]{@{}
  >{\raggedleft\arraybackslash}p{(\linewidth - 4\tabcolsep) * \real{0.3333}}
  >{\raggedright\arraybackslash}p{(\linewidth - 4\tabcolsep) * \real{0.3333}}
  >{\raggedright\arraybackslash}p{(\linewidth - 4\tabcolsep) * \real{0.3333}}@{}}
\caption{Team Roles}\tabularnewline
\toprule\noalign{}
\# & Role & Responsibility \\
\midrule\noalign{}
\endfirsthead
\toprule\noalign{}
\# & Role & Responsibility \\
\midrule\noalign{}
\endhead
\bottomrule\noalign{}
\endlastfoot
1 & Developer & \begin{itemize}
\tightlist
\item
  Write \textbf{quality} code
\item
  Review pull requests
\end{itemize} \\
2 & Designer & \begin{itemize}
\tightlist
\item
  Create mockups
\item
  Maintain style guide
\end{itemize} \\
3 & Manager & \begin{itemize}
\tightlist
\item
  Coordinate team
\item
  Report status
\end{itemize} \\
\end{longtable}

\subsection{Horizontal table}\label{horizontal-table}

\protect\phantomsection\label{tbl:measurements}
\begin{longtable}[]{@{}lll@{}}
\caption{Sensor Readings}\tabularnewline
\toprule\noalign{}
Current (mA) & Voltage (V) & Power (mW) \\
\midrule\noalign{}
\endfirsthead
\toprule\noalign{}
Current (mA) & Voltage (V) & Power (mW) \\
\midrule\noalign{}
\endhead
\bottomrule\noalign{}
\endlastfoot
15 & 3 & 45 \\
16 & 3.1 & 49.6 \\
14 & 2.9 & 40.6 \\
\end{longtable}

\subsection{Table from template}\label{table-from-template}

\protect\phantomsection\label{tbl:project}
\begin{longtable}[]{@{}ll@{}}
\caption{Project Info}\tabularnewline
\toprule\noalign{}
Field & Value \\
\midrule\noalign{}
\endfirsthead
\toprule\noalign{}
Field & Value \\
\midrule\noalign{}
\endhead
\bottomrule\noalign{}
\endlastfoot
Name & pandoc-utils \\
Language & Python, Lua, Bash \\
License & MIT \\
\end{longtable}

\end{document}
